\chapter{Stefan--Boltzmann Law}

\section*{Introduction}

Every object with a temperature above absolute zero emits electromagnetic radiation. The Stefan--Boltzmann law maps the temperature of an object onto the total power it radiates, with no dependence on the details of the emission process. It is the macroscopic expression of the microscopic quantum statistics of photons, and it explains why hot objects glow---and why the color of that glow shifts with temperature.


\begin{equation}
j = \sigma T^4
\end{equation}

\section*{Fundamental Equations Used}

The derivation starts from Planck's spectral energy density of black-body radiation, integrated over all frequencies.

\section*{Assumptions}

\textbf{Assumptions:} The object is a perfect black body (absorptivity = emissivity = 1); thermal equilibrium; the radiation fills a cavity or the object is in vacuum.


\textbf{Limitations:} Real objects have emissivity $\epsilon < 1$, so the actual radiated power is $j = \epsilon\sigma T^4$; the law does not give the spectral distribution (Planck's law does).


\section*{Mathematical Derivation}

\textbf{Step 1: Write Planck's spectral radiance.}

The spectral energy density of black-body radiation at temperature $T$ is
\begin{equation}
u(\nu,T) = \frac{8\pi h\nu^3}{c^3}\,\frac{1}{e^{h\nu/k_BT} - 1},
\end{equation}
where $8\pi\nu^2/c^3$ is the density of photon modes per unit volume per unit frequency, and $h\nu/(e^{h\nu/k_BT}-1)$ is the mean energy per mode (Planck distribution).

\textbf{Step 2: Compute the total energy density.}

Integrating over all frequencies:
\begin{equation}
u(T) = \int_0^\infty u(\nu,T)\,d\nu = \frac{8\pi h}{c^3}\int_0^\infty \frac{\nu^3}{e^{h\nu/k_BT} - 1}\,d\nu.
\end{equation}
Substitute $x = h\nu/(k_BT)$, so $\nu = k_BTx/h$ and $d\nu = (k_BT/h)\,dx$:
\begin{equation}
u(T) = \frac{8\pi h}{c^3}\left(\frac{k_BT}{h}\right)^4\int_0^\infty \frac{x^3}{e^x - 1}\,dx.
\end{equation}

\textbf{Step 3: Evaluate the standard integral.}

The integral $\int_0^\infty x^3/(e^x - 1)\,dx = \pi^4/15$ (a standard result from Bose-Einstein statistics). Thus:
\begin{equation}
u(T) = \frac{8\pi^5 k_B^4}{15\,h^3 c^3}\,T^4.
\end{equation}

\textbf{Step 4: Relate energy density to radiated flux.}

The intensity (power per unit area) radiated by a black body is related to the energy density by $j = uc/4$. The factor $1/4$ arises because the radiation is isotropic: the flux through a surface is $u\langle v\cos\theta\rangle/4$ where $\langle v\rangle = c$ and $\langle\cos\theta\rangle = 1/2$ over the hemisphere. Thus:
\begin{equation}
j = \frac{uc}{4} = \frac{\pi^5 k_B^4}{60\,h^3 c^2}\,T^4.
\end{equation}

\textbf{Step 5: Identify the Stefan-Boltzmann constant.}

Defining $\sigma = \pi^5 k_B^4/(60\,h^3 c^2)$:
\begin{equation}
\boxed{j = \sigma\,T^4, \qquad \sigma = \frac{\pi^5 k_B^4}{60\,h^3 c^2} = 5.670 \times 10^{-8}\;\text{W}\,\text{m}^{-2}\,\text{K}^{-4}.}
\end{equation}
The $T^4$ dependence follows from dimensional analysis: the only combination of $k_B$, $h$, $c$ with units of power per area times $T^{-4}$ gives $\sigma$. The $T^4$ power law is a direct consequence of the Bose-Einstein statistics of photons.

\section*{Characteristics of the Equation}

\begin{itemize}
    \item The $T^4$ dependence is extremely strong: doubling the temperature increases the radiated power by a factor of 16.
    \item The Stefan-Boltzmann constant $\sigma$ can be derived from Planck's radiation law by integrating over all wavelengths.
    \item Joseph Stefan derived the $T^4$ law empirically in 1879; Ludwig Boltzmann derived it theoretically from thermodynamics in 1884.
    \item The total power from an object of area $A$ and emissivity $\epsilon$ is $P = \epsilon\sigma A T^4$.
\end{itemize}
\section*{Solution Methods}

\begin{itemize}
    \item \textbf{From Planck's law:} Integrate $B_\lambda(T)$ over all wavelengths: $\sigma = \int_0^\infty \pi B_\lambda(T)\,d\lambda / T^4$.
    \item \textbf{Thermodynamic derivation:} Use the radiation pressure $P = u/3$ and the first law of thermodynamics to derive the $T^4$ scaling.
    \item \textbf{Dimensional analysis:} The only combination of $k$, $h$, $c$, and $T$ with units of power per area is $\sigma T^4$.
\end{itemize}
\section*{Verifications}

\begin{itemize}
    \item The measured solar constant (1361 W/m$^2$) combined with the Stefan-Boltzmann law gives the Sun's effective temperature as approximately 5778 K, consistent with spectroscopic measurements.
    \item Laboratory black-body cavities confirm the $T^4$ dependence to high precision.
    \item The Cosmic Microwave Background spectrum matches a black body at $T = 2.725$ K, with the total power given by the Stefan-Boltzmann law.
\end{itemize}
\section*{Applications}

\begin{itemize}
    \item Stellar astrophysics: determining stellar luminosities and temperatures from observed flux.
    \item Climate science: Earth's energy balance and the greenhouse effect.
    \item Infrared thermography and pyrometry for non-contact temperature measurement.
    \item Design of thermal protection systems for spacecraft and re-entry vehicles.
\end{itemize}
\section*{What Richard Feynman Would Have Asked}

\subsection*{Plain-English Translation}
\textit{``Why does hot stuff glow, and why does the Stefan-Boltzmann law say the power goes as $T^4$ rather than some other power?''}

Hot stuff glows because every object emits electromagnetic radiation due to the thermal motion of its charged particles. The $T^4$ dependence comes from two effects working together: as temperature increases, each frequency mode of the radiation contains more energy (proportional to $T$), and more frequency modes become active (proportional to $T^3$). The product gives $T^4$. This is why a modest increase in temperature causes a dramatic increase in brightness---the fourth power is very steep.

\subsection*{Localized Mechanics}
\textit{``What is happening at a single point on the surface of a hot object that causes it to radiate? Is it the atoms vibrating?''}

At any point on the surface, the atoms and their electrons are in constant thermal motion. Accelerating charges emit electromagnetic radiation, and the thermal agitation provides the acceleration. The hotter the surface, the faster the particles move and the more radiation they emit. The Stefan-Boltzmann law is the collective result of countless such emitters, each contributing a small amount of radiation whose frequency spectrum is determined by the temperature. The surface acts as a statistical ensemble of oscillators, each radiating according to quantum mechanics.

\subsection*{Testing Extreme Limits}
\textit{``What happens to the Stefan-Boltzmann law in the extreme limit of very high temperature, like inside a star, or very low temperature, near absolute zero?''}

At very high temperatures, the law still holds as long as the object remains in thermal equilibrium, but new physics enters: at $T > 10^9$ K, electron-positron pair production contributes to the radiation, and the effective number of degrees of freedom changes. At very low temperatures, the $T^4$ law still holds for the radiation field itself, but the coupling between matter and radiation weakens, and the total power becomes extremely small. Near absolute zero, the law gives essentially zero power, consistent with the third law of thermodynamics.

\subsection*{Dimensionality and Geometry}
\textit{``How would the Stefan-Boltzmann law change if we lived in a universe with more than three spatial dimensions?''}

In $d$ spatial dimensions, the density of photon states scales as $T^{d-1}$ instead of $T^3$ (from the surface of the $(d-1)$-sphere in momentum space). Combined with the energy per mode ($\sim kT$), the total radiated power per unit area would scale as $T^{d+1}$ instead of $T^4$. In two dimensions, it would be $T^3$; in four spatial dimensions, $T^5$. The Stefan-Boltzmann constant would also change, absorbing the modified density of states. The qualitative physics---hot objects radiate more---remains the same, but the quantitative scaling changes with dimensionality.

\subsection*{Cross-Disciplinary Connection}
\textit{``Where does the same $T^4$ or power-law scaling appear in other areas of physics?''}

The $T^4$ scaling is specific to photon radiation in three dimensions, but power-law scaling of flux with temperature appears elsewhere: the Debye $T^3$ law for low-temperature specific heat (phonons instead of photons), the $T^2$ scaling of electron specific heat in metals, and the $T^{3/2}$ scaling of magnon contributions in ferromagnets. In each case, the exponent reflects the density of states and the dispersion relation of the relevant excitations. The Stefan-Boltzmann law is the prototype for all these ``power-law thermal response'' phenomena, and the technique of integrating the density of states over energy is the common thread connecting them.
\section*{FAQ}

\textbf{Q: Does a perfect black body really exist?}
A: No real object is a perfect black body, but some come very close. A small hole in a cavity is an excellent approximation, as are certain specially coated surfaces. Real objects have emissivity $\epsilon < 1$, so they radiate less than the Stefan-Boltzmann limit.

\textbf{Q: Why is the exponent exactly 4?}
A: The $T^4$ arises from the integration of the Planck spectrum over all wavelengths. Each wavelength contributes proportionally to $T$, and the number of modes per unit volume scales as $T^3$ (from the density of photon states). Together, $T \times T^3 = T^4$. Thermodynamically, it follows from the relation between radiation pressure and energy density.

\textbf{Q: Can the Stefan-Boltzmann law be used to measure the temperature of distant objects?}
A: Yes, this is the basis of infrared pyrometry and stellar temperature measurement. If you know the emissivity and measure the total radiated flux, you can solve for $T = (j/(\epsilon\sigma))^{1/4}$. For stars, the black-body approximation is excellent.
\section*{Important Bibliography}

\begin{enumerate}
\item Planck, M., \textit{The Theory of Heat Radiation} (translated by M.\ Masius), 2nd ed. (Dover, 1959).
\item Kittel, C. and Kroemer, H., \textit{Thermal Physics}, 2nd ed. (W.H.\ Freeman, 1980), Chapter 8.
\item Stefan, J., ``\"{U}ber die Beziehung zwischen der W\"{a}rmestrahlung und der Temperatur,'' \textit{Sitzungsber.\ Akad.\ Wiss.\ Wien} \textbf{79}, 391--428 (1879).
\item Rybicki, G.B. and Lightman, A.P., \textit{Radiative Processes in Astrophysics}, 2nd ed. (Wiley-VCH, 1979), Chapter 1.
\end{enumerate}

